\begin{circuitikz}[
    american,
    voltage dir=RP,
    >=latex,
    %>={Latex[length=5mm, width=2mm]},
    alt={Ladder logic diagram bounded by two vertical power rails. The left power rail is highlighted in green. The diagram consists of one active rung. On the rung, a line extends from the left rail through a Greater Than or Equal comparison block labeled GE. Inside the GE block, Source A is set to tag PartCount with a current value of 10, and Source B is set to tag BatchLimit with a value of 10. Following the GE comparison block, the line connects to an output coil labeled OpenDoor, which connects to the right power rail. The entire rung, including the GE comparison block and the OpenDoor coil, is highlighted in green through to the right power rail because the comparison evaluates to true.}
    ]

    % Define the number of rungs in the diagram
    \def\numRungs{1}

    % Calculate the length of the vertical power rails dynamically based on number of rungs
    \pgfmathsetmacro{\railLength}{-2 * \numRungs}

    %======================================================================================================
    % Component Grid
    % Spacing = 3x, 2y
    %======================================================================================================
    \foreach \i in {0,...,5} {
        \foreach \j in {0,...,20} { % Change to 0,...,19 if you need 20 individual Y points down to -40
        
        % Calculate spatial positions
        \pgfmathsetmacro{\px}{\i * 3}
        \pgfmathsetmacro{\py}{-1 - (\j * 2)}
        
        % Name coordinate using loop indices: (C-column-row) e.g., (C-0-0), (C-5-10)
        \coordinate (C-\i-\j) at (\px, \py);
        
        % OPTIONAL: Uncomment the 2 lines below to display grid labels while building
        % \fill[red] (C-\i-\j) circle (1.5pt);
        %\node[anchor=south west, scale=0.5, gray] at (C-\i-\j) {(\i,\j)};
        }
    }

    %======================================================================================================
    % Foreground Layer
    % Only contains the vertical power rails
    % Length of the rails is automatic, calculated based on the number of rungs defined above
    %======================================================================================================
    \begin{pgfonlayer}{ft}

        %======================================================================================================
        % Scope sets options for both lines - BakerBlack, very thick
        %======================================================================================================
        \begin{scope}[
            draw=BakerBlack,
            text=BakerBlack,
            color=BakerBlack,
            font=\small,
            ultra thick
        ]
        
            \draw (0,0) -- (0,\railLength);

            \draw (15,0) -- (15,\railLength);

        \end{scope}
    \end{pgfonlayer}

    %======================================================================================================
    % Main Layer
    % All actual logic goes here
    %======================================================================================================
    \begin{pgfonlayer}{main}
        
        %======================================================================================================
        % Scope sets options for both lines - BakerBlack for draw and text
        %======================================================================================================
        \begin{scope}[
            draw=BakerBlack,
            text=BakerBlack,
            color=BakerBlack,
            font=\small
        ]

            %=================================================================
            % Rung 1
            %=================================================================
            \draw (C-0-0) -- ++(0.5,0) pic (GE) {SrcAB};
            \draw (GE-OUT) -- (C-4-0) to[esource, l={OpenDoor}] (C-5-0);

            % GE box labels
            \node[anchor=south] at (GE-TOP) {GE};
            \node[anchor=west] at (GE-SRCA) {PartCount};
            \node[anchor=west] at (GE-AVAL) {10};
            \node[anchor=west] at (GE-SRCB) {BatchLimit};
            \node[anchor=west] at (GE-BVAL) {10};

        \end{scope}
    \end{pgfonlayer}

    %=================================================================
    % Background layer
    %=================================================================
    \begin{pgfonlayer}{bg}
        
        %==================================================================
        % Highlight Scope
        % This layer exists for highlighting true logic, the scope sets the
        % highlight options without having to repeat them every \draw
        % draw=UpNorthGreen, line width=10pt, opacity=0.4
        %==================================================================
        \begin{scope}[
            draw=UpNorthGreen,
            line width=10pt,
            opacity=0.4
        ]

            %==================================================================
            % Left vertical rung - leave this alone, it should ALWAYS be green
            % It automatically draws itself to the length defined at the top
            %==================================================================
            \draw (0,0) -- (0,\railLength);
        
            %=================================================================
            % Rung 1
            %=================================================================
            \draw (C-0-0) -- (GE-IN);

            \draw (GE-OUT) -- (C-5-0);

        \end{scope}
    \end{pgfonlayer}
    
\end{circuitikz}